\documentclass[tikz,border=10pt]{standalone}
\usetikzlibrary{calc}

\begin{document}
\begin{tikzpicture}
    % Define coordinates
    % Base BC
    \coordinate (B) at (0,0);
    \coordinate (C) at (10,0);
    
    % Point E such that angle BEC = 90 degrees.
    % Placed at approx (2.5, 4.33) to satisfy the right angle with B and C.
    \coordinate (E) at (2.5, 4.3301);
    
    % Point A such that angle BAC = 90 degrees.
    % Placed symmetrically at approx (7.5, 4.33).
    \coordinate (A) at (7.5, 4.3301);
    
    % Point D is the intersection of AB and EC.
    % Calculated based on the lines.
    \coordinate (D) at (5, 2.8867);

    % Draw edges
    \draw (B) -- (C);
    \draw (B) -- (E);
    \draw (E) -- (C);
    \draw (A) -- (C);
    \draw (A) -- (B);
    
    % Draw right-angle mark at E (between BE and EC)
    \coordinate (E1) at ($(E)!0.3cm!(B)$);
    \coordinate (E2) at ($(E)!0.3cm!(C)$);
    \draw (E1) -- ($(E1)+(E2)-(E)$) -- (E2);
    
    % Draw right-angle mark at A (between AB and AC)
    \coordinate (A1) at ($(A)!0.3cm!(B)$);
    \coordinate (A2) at ($(A)!0.3cm!(C)$);
    \draw (A1) -- ($(A1)+(A2)-(A)$) -- (A2);
    
    % Labels
    \node[above left] at (E) {$E$};
    \node[above right] at (A) {$A$};
    \node[below left] at (B) {$B$};
    \node[below right] at (C) {$C$};
    \node[below] at (D) {$D$};
\end{tikzpicture}
\end{document}
